\section{Quantitative counting and coefficient cancellation}
\label{sec:quantitativecounting}

The preceding construction raises a quantitative question: can explicit
crossing resolutions give a counting remainder with a specified exponent
and geometric constants? We first give an elementary counting benchmark.
We then prove a systole-dependent gap and a uniform recovery estimate,
and show on a fixed surface why their exponential separation scale is
unavoidable for estimates of individual orbital terms. Finally, a finite
support coefficient and a smoothing criterion isolate the additional
estimate needed for connected-curve counting.

\subsection{Effective counting and the meaning of an explicit saving}
\label{subsec:countknown}

Continue with a closed hyperbolic surface $X$ of genus $g\geq2$, and
put $d=6g-6$.  For a fixed nonzero rational multicurve $\gamma_0$, write
\[
 N_{X,\gamma_0}(R)=
 \#\{\gamma\in\operatorname{Mod}(S)\gamma_0:
                         \ell_X(\gamma)\leq R\}.
\]
A remainder $O(R^{d-\kappa})$ with $\kappa>0$ is called a
power saving because it is smaller than the $R^d$ main term
by a factor $R^{-\kappa}$. The subscript on $O_{X,\gamma_0}$
allows its constant to depend on the metric and curve type.
An existence theorem for a positive $\kappa$ and an explicit
numerical value for it are different levels of quantitative information.

Eskin--Mirzakhani--Mohammadi prove that there exists
$\kappa(g)>0$ such that
\begin{equation}
 N_{X,\gamma_0}(R)
 =c_{X,\gamma_0}R^d
   +O_{X,\gamma_0}(R^{d-\kappa(g)}).
 \label{eq:countEMM}
\end{equation}
Their theorem also allows variable negative curvature; the hyperbolic
case suffices here~\cite[Theorem 1.1]{EMMEffective}.  The theorem
states the existence of the exponent rather than a numerical value.
Its proof uses exponential mixing of the Teichm\"uller flow.  A
specified exponent and certified constants for a given surface would
therefore be additional quantitative information.  This distinction
does not diminish the proved power saving or establish that its
constants are intrinsically noncomputable.

The coefficient distributions offer explicit local formulas, but
their present estimates do not improve~\eqref{eq:countEMM}.
The task is to control a growing collection of connected supports.
The qualitative leading term, treated in
\cite{MirzSimple,RivinSimpleAsymptotics}, and an effective remainder
are separate questions.  It is useful first to record how much
quantitative control follows from ordinary lattice geometry alone.

\subsection{A constructive lattice benchmark for all multicurves}
\label{subsec:countlattice}

A measured lamination is a closed union of disjoint geodesic leaves
with a compatible transverse measure. Weighted disjoint closed
curves are its integral examples, and geometric length extends
continuously and homogeneously to this space. A train track is a
finite branched diagram carrying such leaves; nonnegative branch
weights satisfy switch equations expressing conservation at each
junction. These equations provide linear coordinate cones and
integral lattices. Thurston measure is the compatible Lebesgue
measure in those coordinates, with its lattice normalization fixed.

Let $\mathcal{ML}_{\ZZ}(S)$ denote integral measured laminations,
equivalently essential simple multicurves with positive integer
weights.  These include disconnected supports.  Normalize Thurston
measure by the actual integral-lamination lattice in each chart, and
set
\[
 M_X(R)=\#\{C\in\mathcal{ML}_{\ZZ}(S)\setminus\{0\}:
                  \ell_X(C)\leq R\},\qquad
 B(X)=\mu_{\mathrm{Th}}\{\lambda:\ell_X(\lambda)\leq1\}.
\]

\begin{proposition}[Total multicurve error]
\label{prop:countall}
For the fixed surface $X$,
\begin{equation}
 \boxed{M_X(R)=B(X)R^d+O_X(R^{d-1}).}
 \label{eq:countall}
\end{equation}
The error can be bounded using finitely many lattice covolumes,
fundamental-cell radii, convex-body inradii and volumes, and
lower-dimensional face bounds.
\end{proposition}

\begin{proof}
Here is the quantitative lattice estimate needed for the argument.
Let $\Lambda\subset\RR^d$ have covolume $v$ and a fundamental cell
$F\subset B(0,\rho)$.  Suppose the convex compact body $K$ has volume
$V$ and contains $B(x_0,r)$, where $r>0$.  Put $a=\rho/r$.  Then
\begin{equation}
 \frac Vv(R-a)_+^d
 \leq\#(\Lambda\cap RK)
 \leq\frac Vv(R+a)^d.
 \label{eq:countsandwich}
\end{equation}
Indeed, up to null sets the lattice cells tile Euclidean space, so
\[
 (RK)\ominus B_\rho
 \ \subset\!
 \bigcup_{p\in\Lambda\cap RK}(p+F)
 \ \subset RK+B_\rho.
\]
Here $A\ominus B_\rho=\{x:x+B_\rho\subset A\}$.
Writing $K_0=K-x_0$, we have $B_r\subset K_0$ and
\[
 RK+B_\rho\subset Rx_0+(R+a)K_0.
\]
For $R\geq a$, also
$Rx_0+(R-a)K_0\subset(RK)\ominus B_\rho$.
Taking volumes proves~\eqref{eq:countsandwich}; the lower bound is
zero when $R<a$.  Notice that the homotheties are centered at
$Rx_0$, since the origin can be a boundary point of $K$.
For $R\geq a$ this gives the explicit bound
\begin{equation}
 \left|\#(\Lambda\cap RK)-\frac VvR^d\right|
 \leq d\,2^{d-1}\frac Vv\frac\rho r\,R^{d-1}.
 \label{eq:countcellerror}
\end{equation}
The upper increment $(R+a)^d-R^d$ dominates the lower decrement.

The scale of the error is familiar from a square. For
$K=[0,1]^2$ and $\Lambda=\mathbb Z^2$, the exact count in
$RK$ is $(\lfloor R\rfloor+1)^2=R^2+O(R)$.
Interior lattice cells account for area; the uncertainty lies in
cells meeting the boundary, whose number has one lower power of
$R$. The inradius and cell-radius argument above supplies this
principle for a general convex body, without assuming a smooth boundary.
The covolume $v$ is the volume of one lattice fundamental cell.

To apply this estimate, use a finite rational polyhedral cone
subdivision subordinate to train-track charts, with their actual
integral weight lattices. Length is convex in each train-track cone,
by the surgery argument in \cite[Theorem 5.1]{EMMEffective}.
The finite polyhedral handling of chart overlaps is given in
\cite[Lemma 5.1 and the end of the proof of Theorem 1.1 in
Section 8]{EMMEffective}; a further finite subdivision makes the
chambers convex and their overlaps faces. In the coordinates of
each chamber $Q_i$, the sets
$K_i=Q_i\cap\{\ell_X\leq1\}$ are compact convex bodies.
Indeed, continuity and positivity of length on the projectivized
cone give boundedness, and scaling an interior point of a
full-dimensional chamber to length less than one gives an interior
ball. Lower-dimensional chambers are handled as faces. Apply
\eqref{eq:countsandwich} in every full-dimensional chamber with its
own integral lattice.  Their normalized volumes sum to $B(X)$.
In particular, a coordinate parity condition must be absorbed into
the lattice covolume, not silently discarded.

Shared faces can be assigned to one chamber.  Alternatively, count
closed chambers and bound the overcount.  If a face has dimension
$1\leq k\leq d-1$, its unit-length part lies in a ball of radius $A_F$, and
its induced lattice has shortest nonzero vector of length $\delta_F$,
disjoint-ball packing gives
\[
 \#(\Lambda_F\cap RK_F)
 \leq\left(1+\frac{2RA_F}{\delta_F}\right)^k.
\]
Zero-dimensional faces contribute at most one point each.  The
finite sum of these bounds is $O_X(R^{d-1})$.
Removing the zero multicurve changes the answer by one.
\end{proof}

This refines the elementary total-multicurve lattice asymptotic in
\cite[Proposition 3.1]{MirzSimple}; it is a benchmark for the new
machinery.  Formula~\eqref{eq:countcellerror} is explicit once bounds
on its geometric inputs are certified.  It does not supply numerical
values for an unspecified surface.  Fixed train-track representatives
give linear upper bounds for length, and coordinate length estimates
such as~\cite{LuoStongLengths} help certify the required geometric
bounds.  Hyperbolic length itself need not be piecewise linear.
For example, choose an interior rational weight vector in a chamber
and scale it until a certified linear upper bound for length is less
than one.  A small coordinate ball then remains both in that chamber
and below the upper bound, producing an inradius certificate.
For an outer-radius certificate one needs a lower bound for length
on the chamber's unit coordinate sphere.  A finite filling collection
of curves, together with collar inequalities and their
piecewise-linear intersection functions, supplies such a bound.
The constants can deteriorate when the surface degenerates; no
uniformity over the entire moduli space is being asserted.
On a fixed compact subset of Teichm\"uller space the same finite
coordinate construction and uniform length comparisons give uniform
versions of the bounds.

\subsection{Primitivity, connectivity, and the torus benchmark}
\label{subsec:countprimitive}

Write $C=\sum_jm_j\alpha_j$ with distinct simple supports.  Call $C$
primitive as an integral lamination if $\gcd_jm_j=1$.
Every nonzero integral multicurve is uniquely $kC_0$ with $C_0$
primitive.  This intrinsic definition avoids coordinate conventions:
division by $k$ is allowed precisely when the resulting lamination
remains integral.  Let $P_X(R)$ count these primitive laminations.
M\"obius inversion gives
\[
 P_X(R)=\sum_{k\leq R/\operatorname{sys}(X)}\mu(k)M_X(R/k).
\]
Since $d>2$, the error series $\sum k^{-(d-1)}$ converges, while
the truncated leading-term tail is $O_X(R)$.  Therefore
\begin{equation}
 P_X(R)=\frac{B(X)}{\zeta(d)}R^d+O_X(R^{d-1}).
 \label{eq:countprimitive}
\end{equation}
Here $\zeta(d)=\sum_{k\ge1}k^{-d}$, and the arithmetic
M\"obius function is the one defined in Section~\ref{sec:coefficientdistributions}.
For example, $2\alpha+3\beta$ is primitive when $\alpha,\beta$
are distinct disjoint simple supports, since $\gcd(2,3)=1$.
It still has two supports. Dividing out a common integer weight
and requiring connected support answer different questions.

A union $\alpha+\beta$ of distinct disjoint simple curves is primitive
but disconnected.  Thus~\eqref{eq:countprimitive} does not count
connected simple curves in higher genus.

For a once-punctured torus every essential multicurve has one
distinct support.  The McShane--Rivin norm on $H_1(T;\RR)$ identifies
simple curves with primitive lattice vectors modulo sign.  If $B_X$
is its full centrally symmetric unit ball, our unoriented convention
gives
\begin{equation}
 N_X(R)=\frac{\operatorname{area}(B_X)}{2\zeta(2)}R^2
                  +O_X(R\log R).
 \label{eq:counttorus}
\end{equation}
This is the norm-and-M\"obius-inversion argument of
\cite[Theorem 4.1]{McShaneRivinTori}; the factor $1/2$ identifies
opposite homology vectors.  The logarithm comes from
$\sum_{k\leq O_X(R)}k^{-1}$.  This independent torus benchmark does
not extend the closed-surface trace formula across a cusp.


\subsection{A fixed-surface obstruction to uniform first-atom separation}
\label{subsec:countnearcollision}

The exponentially small remainder for one fixed simple curve does not
remain small uniformly over curves whose length tends to infinity.
There is an explicit obstruction on a fixed hyperbolic surface, with
no degeneration of its metric.

\begin{proposition}[An exponentially close pants contribution]
\label{prop:countnearcollision}
On every closed hyperbolic surface $X$ of genus at least two there
are fixed distinct simple curves $a,b$, and simple curves $c_n$
with $R_n=\ell_X(c_n)\to\infty$, such that $a,b,c_n$ bound
embedded geodesic pairs of pants. Let $\gamma_n$ be their
figure-eight curves with
$T_{\gamma_n}=T_aT_b-T_{c_n}$, and set
\[
 A=\cosh\frac{\ell_X(a)}2\cosh\frac{\ell_X(b)}2,
 \qquad L_n=\ell_X(\gamma_n).
\]
Then
\begin{align}
 L_n-R_n&=8A e^{-R_n/2}+O(e^{-R_n}),\notag\\
 L_n^2-R_n^2&=16A R_n e^{-R_n/2}
                    +O(R_n e^{-R_n}).
 \label{eq:countpantsgap}
\end{align}
The coefficient of $T_{c_n}$ in $T_{\gamma_n}$ is $-1$.
\end{proposition}

\begin{proof}
Choose distinct disjoint nonseparating curves $a,b$ whose complement
after cutting is connected, of genus $g-2$ with four boundary
components. In this cut surface let $c$ bound a pair of pants with
one copy of $a$ and one copy of $b$. Choose an essential interior
curve intersecting $c$. A Dehn twist about that curve fixes the boundary
and gives infinitely many distinct curves $c_n$, each bounding a
pair of pants with $a,b$. Straightening preserves this configuration.
Finiteness of the closed-geodesic length set below a bound implies
$R_n\to\infty$.

Rivin's pants identity~\cite[Theorem 1.3]{RivinFigures} gives
\[
 \cosh(L_n/2)=\cosh(R_n/2)+2A.
\]
Expanding its inverse for $R_n\to\infty$ gives
\eqref{eq:countpantsgap}. The coefficient assertion is the ordinary
trace identity in the displayed pants variables.
\end{proof}

Normalize the $c_n$ heat coefficient by its first-atom contribution.
The single primitive orbital term supplied by $\gamma_n$ becomes
\begin{equation}
 -r_n\exp\!\left(-\frac{L_n^2-R_n^2}{4t}\right),
 \qquad
 r_n=\frac{L_n}{R_n}
       \frac{\sinh(R_n/2)}{\sinh(L_n/2)}
     =1+O(e^{-R_n/2}).
 \label{eq:countnormalizedpants}
\end{equation}
For $t=R_n^{-p}$, with any fixed $p>0$, this tends to $-1$.
Thus the absolute-valued orbital remainder contains an order-one
term at every polynomial time scale. Suppressing this one term
below $R_n^{-M}$ requires
\begin{equation}
 t\leq\left(\frac{4A}{M}+o(1)\right)
       \frac{R_n e^{-R_n/2}}{\log R_n}.
 \label{eq:countnecessarytime}
\end{equation}
This is a statement about an individual orbital contribution and
the absolute-valued remainder. It does not exclude cancellations
with other curves, including other curves of exactly the same
length, and does not by itself identify a nonzero net atom of
$\sigma_{c_n}^X$.

Such cancellation is precisely what the coefficient approach can
organize. By~\eqref{eq:coefffigurecancellation},
$\sigma_{c_n}^X+\sigma_{a+b}^X$ cancels the entire chosen
$\gamma_n$ family, including all iterates, exactly. A useful
counting construction should exploit these relations before taking
absolute values. The proposition explains why the cancellation is
more than a small correction to an already uniform first-atom
estimate.

\subsection{A constructive uniform isolation estimate}
\label{subsec:countuniformisolation}

The same bracelet estimate that bounded the coefficients also
provides an explicit lower bound on their orbital separation.
Write $\varepsilon=\operatorname{sys}(X)>0$, and put
\begin{equation}
 c_\varepsilon=\min\left\{1,
       2\cosh(3\varepsilon/2)-2\cosh(\varepsilon/2)\right\}>0.
 \label{eq:counttracegapconstant}
\end{equation}

\begin{lemma}[A quantitative trace gap]
\label{lem:counttracegap}
Let $\alpha$ be simple, $x=\ell_X(\alpha)$, and let
$\gamma=\delta^k$ be an orbital label contributing to
$\sigma_\alpha^X$, other than $(\alpha,1)$. If $L=\ell_X(\gamma)$,
then
\begin{equation}
 2\cosh(L/2)-2\cosh(x/2)\geq c_\varepsilon.
 \label{eq:counttracegap}
\end{equation}
Consequently, for $x\leq R$,
\begin{equation}
 L^2-x^2\geq\Delta_\varepsilon(R),\qquad
 \Delta_\varepsilon(R)=
 4\arcosh^2\!\left(\cosh(R/2)+\frac{c_\varepsilon}{2}\right)-R^2
 \sim4c_\varepsilon R e^{-R/2}.
 \label{eq:countactiongap}
\end{equation}
\end{lemma}

\begin{proof}
By~\eqref{eq:coeffbraceletpositive}, the positive Fuchsian value
$2\cosh(L/2)$ is an integer positive sum of bracelet values.
At least one contributing bracelet has an $\alpha$ coefficient.
It therefore has an odd winding $n\geq1$ around $\alpha$;
any additional support components must contribute constant
Chebyshev terms and have even windings at least two.

If $n\geq3$, that bracelet exceeds $2\cosh(x/2)$ by at least
$2\cosh(3\varepsilon/2)-2\cosh(\varepsilon/2)$.
If $n=1$ and additional support occurs, its value is at least
$2\cosh(x/2)\,2\cosh\varepsilon$, giving a difference greater
than two. The remaining case is the bracelet $\alpha$ itself.
Unless it is the only term with coefficient of absolute value one,
integrality gives an additional positive value of at least one.
If it were that sole term, the Fuchsian lengths would satisfy
$L=x$, contrary to the strict length-ordering lemma for the
orbital labels under consideration. This proves
\eqref{eq:counttracegap}.

Set $f(x)=2\arcosh(\cosh(x/2)+c_\varepsilon/2)$.
The trace gap gives $L\geq f(x)>x$. Moreover,
\[
 \frac{\dd}{\dd x}\bigl(f(x)^2-x^2\bigr)
 =2\left[f(x)\frac{\sinh(x/2)}{\sinh(f(x)/2)}-x\right]<0,
\]
because $y/\sinh(y/2)$ decreases for $y>0$.
This proves the uniform action bound; its asymptotic follows by
expanding $f(R)-R=2c_\varepsilon e^{-R/2}+O(e^{-R})$.
\end{proof}

Here and below ``orbital gap'' refers to the individual summands
before cancellation at coincident lengths. This convention both
makes the estimate stronger and avoids assuming that a later
summand survives in the signed measure.

For explicit constants let $D=\operatorname{diam}(X)$ and define
\begin{align}
 Q_X&=\frac{e^{2D}}{(1-e^{-\varepsilon/2})^2},\notag\\
 A_\varepsilon&=\frac{1+e^{-\varepsilon}}{(1-e^{-\varepsilon})^2}
       \geq\sup_{\substack{n\geq1\\n\ {\rm odd}}}
             n e^{-(n-1)\varepsilon/2},\notag\\
 C_X&=\frac{Q_X A_\varepsilon\coth(\varepsilon/2)e^2}
               {(1-e^{-1})^2}\left(1+\frac2\varepsilon\right).
 \label{eq:countuniformconstants}
\end{align}
For a simple $\alpha$ let $b_\alpha$ be the smaller of
$\ell_\alpha+1$ and its first later contributing orbital length.
The latter exists: the orbital label $(\alpha,3)$ contributes
coefficient $-3$, irrespective of cancellation with other labels
at the same length.
For $R\geq\varepsilon$ put
\[
 \eta_X(R)=\min_{\substack{\alpha\ {\rm simple}\\\ell_\alpha\leq R}}
             (b_\alpha^2-\ell_\alpha^2),\qquad
 \underline\eta_X(R)=
             \min\{2\varepsilon+1,\Delta_\varepsilon(R)\}.
\]
Only finitely many labels enter the first minimum, and
Lemma~\ref{lem:counttracegap} gives
$\eta_X(R)\geq\underline\eta_X(R)>0$.

\begin{proposition}[Uniform recovery with explicit constants]
\label{prop:countuniformisolation}
For $0<t\leq\varepsilon/4$,
\begin{equation}
 \boxed{
 \sup_{\substack{\alpha\ {\rm simple}\\\ell_\alpha\leq R}}
 \left|
   \frac{\sinh(\ell_\alpha/2)}{\ell_\alpha}
   e^{\ell_\alpha^2/(4t)}G_\alpha^X(t)-1
 \right|
 \leq C_X\exp\!\left(R-\frac{\eta_X(R)}{4t}\right).
 }
 \label{eq:countuniformisolation}
\end{equation}
The explicit $\underline\eta_X(R)$ may replace $\eta_X(R)$.
\end{proposition}

\begin{proof}
Orbit packing gives a bound $N_{\rm all}(U)\leq Q_Xe^U$
for primitive/iterate labels of length at most $U$.
Indeed, choose a representative whose axis comes within $D$ of
a fixed lift of a basepoint. It moves that point at most $U+2D$.
The orbit points have separation at least $\varepsilon$, and
comparison of disjoint disks of radius $\varepsilon/2$ with the
disk of radius $U+2D+\varepsilon/2$ gives the stated constant.

Set $x=\ell_\alpha$ and $w_\alpha=x/\sinh(x/2)$.
The sharper form of the fixed-coefficient bound is
\[
 \frac{K_{\alpha,X}}{w_\alpha}
 \leq\frac{A_\varepsilon}{2x}.
\]
In fact,
$\sinh(x/2)/\cosh(nx/2)\leq e^{-(n-1)x/2}$, and the sum
over odd $n=2j+1$ of
$(2j+1)e^{-j\varepsilon}$ equals $A_\varepsilon$.
Thus the absolute mass of a length-$L$ orbital term, divided by
$w_\alpha$, is at most
$A_\varepsilon\coth(\varepsilon/2)L/x$.
Sum the tail in unit shells starting at $b=b_\alpha$.
Since $b\geq\varepsilon$ and $t\leq\varepsilon/4$, its normalized
absolute value is at most the following expression: the shell
count contributes $e^j$, whereas
$((b+j)^2-b^2)/(4t)\geq2j$ supplies at least $e^{-2j}$.
\[
 \frac{Q_X A_\varepsilon\coth(\varepsilon/2)}{x}
 e^{b+1-(b^2-x^2)/(4t)}
 \sum_{j\geq0}(b+j+1)e^{-j}
 \leq
 \frac{Q_X A_\varepsilon\coth(\varepsilon/2)}{(1-e^{-1})^2}
 \frac{b+1}{x}e^{b+1-(b^2-x^2)/(4t)}.
\]
Use $b\leq x+1\leq R+1$ and $(b+1)/x\leq1+2/\varepsilon$.
This proves~\eqref{eq:countuniformisolation}.
\end{proof}

For $R>1$ and $M>0$, a sufficient condition for relative error
at most $R^{-M}$ is
\begin{equation}
 t\leq\min\left\{\frac\varepsilon4,
 \frac{\underline\eta_X(R)}
      {4\bigl(R+M\log R+\log\max\{1,C_X\}\bigr)}\right\}.
 \label{eq:countcertifiedtime}
\end{equation}
Every displayed constant is determined by geometric data, and
certified lower bounds for the systole and upper bounds for the
diameter suffice for a conservative version. A sharper $\eta_X(R)$
can also be found from the finite cutoff $R+1$ when length
comparisons and the metric are supplied effectively. The guaranteed
time scale in~\eqref{eq:countcertifiedtime} is exponential in $-R/2$.
Together with Proposition~\ref{prop:countnearcollision}, this
identifies the exponential separation scale constructively.

This is a quantitative recovery theorem, not a counting asymptotic.
Summing its normalized expression over $\ell_\alpha\leq R$
already presupposes the target set and its lengths. A counting
argument must construct an aggregate observable, perform the
necessary signed cancellations, and estimate its large-length
behavior without assuming that list.

\subsection{A finite coefficient that selects one support}
\label{subsec:countsupport}

There is a useful way to state the remaining selection problem without
an infinite-dimensional coefficient operation. For an integral
multicurve $C=\sum_{i=1}^r m_i\alpha_i$, with distinct supports
$\alpha_i$, put $r(C)=r$ and $r(0)=0$. Define
\begin{equation}
 \mathcal P_X(s,z)=\sum_{C\in\mathcal{ML}_{\ZZ}(S)}
       z^{r(C)}e^{-s\ell_X(C)},\qquad \Re s>0.
 \label{eq:countsupportpartition}
\end{equation}
The estimate~\eqref{eq:countall} implies absolute
convergence. Since $r(C)\leq3g-3$, this is a polynomial in $z$ of
degree at most $3g-3$, with holomorphic coefficients in $s$.
Here $r(C)$ counts distinct geodesic supports, not the number of
parallel copies used in the trace monomial $T_C$.

Its linear coefficient has an exact interpretation:
\begin{equation}
 F_X(s):=[z]\mathcal P_X(s,z)
 =\sum_{\alpha\ {\rm simple}}\sum_{m\geq1}e^{-sm\ell_\alpha}
 =\sum_{\alpha\ {\rm simple}}\frac{1}{e^{s\ell_\alpha}-1}.
 \label{eq:countonesupport}
\end{equation}
All curves in these sums are unoriented. Divisor inversion gives
\begin{equation}
 \boxed{\quad
 Z_X^{\rm simple}(s):=\sum_{\alpha\ {\rm simple}}e^{-s\ell_\alpha}
 =\sum_{k\geq1}\mu(k)F_X(ks),\qquad \Re s>0.
 \quad}
 \label{eq:countsupportmobius}
\end{equation}
There is no factor $1/k$. Indeed, after expanding the right side,
the coefficient of $e^{-sn\ell_\alpha}$ is
$\sum_{k\mid n}\mu(k)$, which is one for $n=1$ and zero otherwise.
The rearrangement is absolutely convergent: for
$\sigma=\Re s>0$ and $\varepsilon=\operatorname{sys}(X)$,
\[
 \sum_{\alpha,k,m\geq1}e^{-\sigma km\ell_\alpha}
 \leq\frac{Z_X^{\rm simple}(\sigma)}
               {(1-e^{-\sigma\varepsilon})^2}<\infty,
\]
using $km\geq k+m-1$.

For a fixed simple $\alpha$, all $m\alpha$ have one support,
so they contribute $z(e^{-s\ell_\alpha}+e^{-2s\ell_\alpha}+\cdots)$.
For disjoint distinct $\alpha,\beta$, a term $2\alpha+3\beta$
contributes $z^2e^{-s(2\ell_\alpha+3\ell_\beta)}$.
This explains why the linear coefficient selects one support
before M\"obius inversion removes its multiplicity.

Even the support-coefficient extraction is explicit and bounded.
Set $q=3g-2$ and $\omega=e^{2\pi i/q}$. Then
\begin{equation}
 F_X(s)=\frac1q\sum_{j=0}^{q-1}
                 \omega^{-j}\mathcal P_X(s,\omega^j).
 \label{eq:countsupportfourier}
\end{equation}
In genus two the degree is at most three and $q=4$.
For $P(z)=a_0+a_1z+a_2z^2+a_3z^3$, the formula reads
$[z]P=[P(1)-iP(i)-P(-1)+iP(-i)]/4=a_1$.
The four roots of unity cancel all powers except the first.
The same finite averaging explains the error bound below.

Thus an error bounded by $E(s)$ at each of these finitely many
values remains bounded by $E(s)$ after extracting the linear
coefficient. This operation has none of the growing-degree
conditioning of arbitrary trace-coefficient extraction.

For example, suppose a future argument establishes
\[
 \mathcal P_X(s,\omega^j)
 =s^{-d}A_X(\omega^j)+O_X(s^{-d+\eta}),
 \qquad 0<s\leq1,
\]
for every $j$, with a specified $0<\eta<d-1$ and a polynomial
$A_X$ of degree at most $3g-3$. If $a_1=[z]A_X(z)$, then
\begin{equation}
 Z_X^{\rm simple}(s)
 =\frac{a_1}{\zeta(d)}s^{-d}+O_X(s^{-d+\eta}).
 \label{eq:countsupportconditional}
\end{equation}
To check the error, split the sum in~\eqref{eq:countsupportmobius}
at $k=1/s$. The small-$ks$ remainder is bounded by
$s^{-d+\eta}\sum k^{-d+\eta}$; this converges because
$d-\eta>1$. The omitted leading tail and the large-$ks$ tail
are $O_X(s^{-1})$, the latter by exponential decay from the
positive systole. Both are absorbed in the displayed error.

This is a conditional transfer of a remainder, not a proof of
the hypothesized expansion. At $z=1$ ordinary lattice counting
already gives
\[
 \mathcal P_X(s,1)
 =\Gamma(d+1)B(X)s^{-d}+O_X(s^{-d+1})
 \qquad(s\downarrow0).
\]
It does not determine the other values in
\eqref{eq:countsupportfourier}. The support count is a discrete
connectivity statistic, and a general weighted multicurve cannot
be assigned that statistic by a continuous function on all of
measured-lamination space. Knowing its value for each integral
point does not establish its density in expanding regions.

Connectivity itself is computationally accessible: the
Agol--Hass--Thurston interval-pairing algorithm counts components
of a normal curve in time polynomial in the triangulation size
and the logarithm of its normal weight
\cite[Theorem 12 and Corollary 13]{AHTComponents}.
For the one-support selector, first pass from $C$ to its primitive
integral divisor $C_0$: then $r(C)=1$ if and only if $C_0$ has one
component. This division takes place in an integral-lamination
chart before supplying a normal representative to the component
algorithm; admissibility and parity must be respected in the
chosen coordinates. A coordinate gcd alone is not a substitute
for this integral-divisibility test.
Counting the components of $C$ without this division would count
parallel copies separately and would not compute $r(C)$.
The missing quantitative input is an average estimate for this
selector, or a cancellation formula that supplies such an estimate.
Moreover,~\eqref{eq:countsupportpartition} is a target positive
partition function. It has not been identified with an evaluation
of the signed coefficient heat expansion. Defining a linear
operation $T_C\mapsto z^{r(C)}$ on finite trace polynomials alone
does not change an orbital length $\ell_\gamma$ into the resolved
length $\ell_X(C)$, and does not supply that identification.


\subsection{An explicit criterion for passing from smooth to sharp counts}
\label{subsec:countdesmooth}

Proposition~\ref{prop:countall} also bounds the number of connected
simple curves in a shell: for $R\geq2$ and $0\leq\Delta\leq R/2$,
\begin{equation}
 \#\{\alpha\text{ simple}:|\ell_X(\alpha)-R|\leq\Delta\}
       =O_X\bigl((\Delta+1)R^{d-1}\bigr).
 \label{eq:countshell}
\end{equation}
Subtract the total-multicurve estimates at the shell endpoints;
one-sided limits handle endpoint atoms.  The simple curves form a
subset, so positivity gives the stated bound.

\begin{proposition}[Conditional quantitative desmoothing]
\label{prop:countdesmooth}
Let $N_X(R)$ count connected unoriented simple curves.  For
$0<h\leq1/2$, choose smooth $0\leq f_h\leq1$
equal to one on $[0,1-h]$ and zero on
$[1+h,\infty)$.  Suppose, with specified $\eta>0$ and
$m\geq0$, a future argument proves uniformly in $R,h$ that
\begin{equation}
 \sum_{\alpha\text{ simple}}f_h(\ell_X(\alpha)/R)
 =c_XR^d\,d\!\int_0^\infty
             f_h(u)u^{d-1}\,\dd u
   +O_X(h^{-m}R^{d-\eta}).
 \label{eq:countsmoothhypothesis}
\end{equation}
Then
\begin{equation}
 \boxed{N_X(R)=c_XR^d+O_X(R^{d-\delta}),\qquad
        \delta=\min\left\{1,\frac\eta{m+1}\right\}.}
 \label{eq:countsharpconditional}
\end{equation}
\end{proposition}

\begin{proof}
The cutoff differs from the sharp indicator only in a shell of
width $h R$.  Equation~\eqref{eq:countshell} bounds this
loss by $O_X(h R^d+R^{d-1})$.  The normalized integral
in~\eqref{eq:countsmoothhypothesis} differs from one by
$O(h)$.  Hence the total error is
\[
 O_X\bigl(h R^d+R^{d-1}
                   +h^{-m}R^{d-\eta}\bigr).
\]
Choose $h=R^{-\delta}$ for sufficiently large $R$, with $\delta$ as in the statement.
\end{proof}

For example, a hypothetical smooth saving $\eta=1$ with a
second-derivative cost $h^{-2}$ gives $\delta=1/3$.
Choosing $h=R^{-1/3}$ makes both the smoothing loss $hR^d$
and the analytic error $h^{-2}R^{d-1}$ equal in order to
$R^{d-1/3}$. Making the cutoff sharper reduces one error while
increasing the other; balancing them is the reason for the exponent loss.

This criterion identifies the necessary analytic input and its
explicit payoff.  The present coefficient estimates do not establish
\eqref{eq:countsmoothhypothesis}.  In particular, fixed-target
first-atom asymptotics must acquire quantitative control of the
sum over growing connected supports before they imply a counting
remainder.

\subsection{The remaining analytic step}

The two limiting regimes must be distinguished. The first-atom theorem
uses $t\downarrow0$ for $G_\alpha(t)$, or a large positive Laplace
parameter. A positive sum over simple curves detects growing lengths
through $Z_X^{\rm simple}(s)$ as $s\downarrow0$. With the denominator
Gaussian $e^{-L^2/(4t)}$, lengths of order $R$ correspond to $t$ of
order $R^2$. These are different asymptotic problems, and the latter
positive measure has not been constructed by summing the present signed
coefficient measures.

The support polynomial~\eqref{eq:countsupportpartition} specifies a
useful aggregate target. Its finite Fourier extraction is stable, and
its primitive inversion is exact. A positive-axis Laplace remainder
such as~\eqref{eq:countsupportconditional} does not by itself justify
an equally strong sharp counting error; quantitative desmoothing or
suitable complex-parameter estimates are still required.
Proposition~\ref{prop:countdesmooth} gives one sufficient condition
with an explicit loss. The unresolved task is to construct and bound
that smoothed connected-support observable using the coefficient
relations, including their pants cancellations, before taking absolute
values. No improved connected-simple-curve exponent is claimed here.


\clearpage
